\section{Related Work}
\label{sec:related}

\subsection{Small Object Detection in Aerial Imagery}
\label{ssec:related_sod}

Detecting micro targets in aerial and maritime imagery is severely hindered by extreme scale variations, low contrast, and complex backgrounds~\cite{VisDrone,UAVDT,TinyPerson,SeaDronesSee}. Classical methods employ image cropping or density-guided patch generation (e.g., ClusDet~\cite{ClusDet}, DMNet~\cite{DMNet}, UFPMP-Det~\cite{UFPMP}), which often require separate coarse detector networks or incur repeated feature extraction on overlapping tiles. Other works design micro-scale label assignment strategies and geometric distance metrics, including Normalized Gaussian Wasserstein Distance (NWD)~\cite{NWD}, RFLA~\cite{RFLA}, and scale-aware balancing (SSABNet)~\cite{SSABNet}. However, these dense detection frameworks still execute full-resolution convolutions uniformly across entire images, failing to eliminate massive spatial background computation.

\subsection{Selective and Dynamic Computation}
\label{ssec:related_selective}

To reduce arithmetic redundancy, selective computation dynamically routes processing to prospective foreground regions. AutoFocus~\cite{AutoFocus} predicts focus pixels to guide multi-scale processing. QueryDet~\cite{QueryDet} accelerates feature pyramids via cascaded sparse queries, while CEASC~\cite{CEASC} formulates adaptive activation masking for sparse convolutions. Recently, ESOD~\cite{ESOD} introduced an early ObjSeeker module after the shallow stem to slice features into patches (AdaSlicer), bypassing background across subsequent stages. Nonetheless, existing selectors rely strictly on a single semantic objectness score and independent per-cell BCE loss, leaving them vulnerable to semantic feature collapse and irrecoverable tiny-target drops.

\subsection{Spectral Saliency and Frequency Methods}
\label{ssec:related_spectral}

Frequency-domain analysis provides complementary spatial representations. SET~\cite{SET} demonstrated that tiny objects preserve distinct high-frequency structural signatures even when semantic activations fade, using frequency background smoothing to enhance dense feature maps. Unlike SET, which applies dense spectral operations over entire multi-scale pyramids, HESOD repurposes spectral and gradient saliency as an ultra-lightweight, channel-pooled routing signal ($\mathcal{O}(1)$ relative to backbone width) to guide recall-safe spatial selection without dense frequency overhead.
